%----------
%	REPRESENTACIÓN DE LAS MUJERES EN LOS MEDIOS
%----------	

La desigualdad en la representación mediática no es un fenómeno reciente, sino un reflejo histórico de asimetrías sociales profundas que se manifiesta de forma evidente en la cobertura de crisis globales. Un caso paradigmático que ilustra la complejidad de este sesgo editorial es la cobertura internacional del tsunami del Océano Índico en 2004. Tal y como denunció la galardonada periodista de investigación y experta en género Deepa Kandaswamy en Women's eNews (2005) \cite{kandaswamy2005}, la respuesta mediática incurrió en una doble negligencia hacia las mujeres: la hipervisibilización sensacionalista y la posterior invisibilización estructural.

Según documentó Oxfam Internacional \cite{oxfam2005}, en diversas zonas afectadas de Indonesia, India y Sri Lanka fallecieron entre tres y cuatro mujeres por cada hombre. Esta drástica brecha de mortalidad fue el resultado directo de roles sociales fuertemente arraigados: las mujeres fallecieron en mayor proporción por encontrarse en sus hogares cuidando de menores o ancianos, por estar esperando en la costa el regreso de los barcos pesqueros, o por carecer de habilidades de supervivencia reservadas culturalmente a los hombres en esas regiones, como trepar árboles o nadar. Como consecuencia del desastre, las supervivientes se enfrentaron a vulnerabilidades extremas derivadas de esta alteración demográfica, incluyendo agresiones sexuales en campos de desplazados, matrimonios forzados de menores y una falta de atención institucional a la pérdida de sus medios de vida.

A pesar de la gravedad de los datos, la prensa \textit{mainstream} ignoró sistemáticamente el informe de Oxfam. Más revelador aún fue el rechazo dentro de la propia industria periodística: cuando se cuestionó esta falta de cobertura en foros internacionales de prensa, las reclamaciones fueron minimizadas por otros profesionales bajo acusaciones de tener una ``mentalidad de gueto femenino'' o un exceso de ``sensibilidad de género'' \cite{kandaswamy2005}. Este episodio histórico demuestra cómo la falta de una mirada crítica en las redacciones no solo genera narrativas incompletas, sino que encubre negligencias estructurales y perpetúa la revictimización.



Lejos de haberse corregido de forma natural, más de dos décadas después aquel desequilibrio estructural apenas ha cambiado. Según los hallazgos del Proyecto de Monitoreo Global de Medios (\textit{Global Media Monitoring Project}, GMMP por sus siglas en inglés) \cite{unwomen2025gmmp} de 2025, el progreso hacia la igualdad de género en las noticias se encuentra estancado, ya que las mujeres solo representan el 26\% de las personas que aparecen, hablan o sobre las que se lee en las noticias impresas y audiovisuales a nivel mundial. Las mujeres continúan siendo una minoría cuando se trata de figurar como protagonistas de la información, constituyendo únicamente el 24\% de los sujetos de las noticias. La brecha se mantiene igual de grave cuando se buscan voces de autoridad, representando apenas el 23\% de las personas consultadas como expertas o portavoces.

Más recientemente, el proyecto de investigación InfoIA publicado en febrero de 2026 por Clara Sáinz de Baranda y Heidy Ríos \cite{sainzdebaranda2026}, que analiza las informaciones sobre IA en los medios digitales españoles, advierte que este ecosistema informativo consolida un ``techo de cristal mediático'' al ser un espacio fuertemente masculinizado. Los datos de este informe revelan que los hombres protagonizan el 41,60\% de las piezas periodísticas frente a un escaso 5,33\% de protagonismo exclusivamente femenino, y acaparan el 61,14\% de las fuentes de autoridad consultadas, relegando a las mujeres al 21,45\%. 

Ahora bien, ambas caras del problema importan. A la infrarrepresentación cuantitativa se suma un sesgo cualitativo que opera en la propia forma del lenguaje, en cómo se nombra y se describe a las mujeres. Se manifiesta, por ejemplo, en el masculino genérico que invisibiliza la autoría femenina (un equipo liderado por investigadoras reducido a ``un equipo de investigadores'', como en el caso de Bindi), en la denominación de las profesionales por sus relaciones familiares en lugar de por su cargo (``la mujer de'' o ``la hija de'' en vez de su nombre y su función) o en el tratamiento asimétrico que describe a una mujer por su aspecto, su edad o su vestimenta mientras presenta a un hombre por su puesto o sus logros.

Este trabajo aborda esa dimensión lingüística y lo hace en sus dos planos. En el plano cuantitativo, detecta la presencia o ausencia de cada fenómeno en cada pieza, lo que permite medir su prevalencia y validarla frente al criterio experto, por ejemplo en qué proporción de las noticias aparece el masculino genérico. En el plano cualitativo, no se queda en la etiqueta: localiza en el texto el fragmento concreto que motiva el veredicto y aporta una explicación, de modo que la analista pueda auditar cada decisión en lugar de aceptarla a ciegas. Ambos planos se operacionalizan en el Capítulo~\ref{sec:methodology}. Detectar estos fenómenos de forma sistemática ha motivado toda una línea de métodos, que se revisa a continuación.



